\section{Conclusion}\label{sec:conclusion}

We gave \lqCCS a reduction semantics modulo structural congruence and a
labelled transition system, both indexed by the choices that tags make
explicit, and proved that they agree: by the Harmony lemma
(\cref{lem:harmony}), a process $\conf{\rho}{P}$ reduces with a given choice
exactly when the composition $\rho \lpar P$ of its memory with its program
performs a silent transition with that choice, up to structural congruence.
In the labelled semantics the quantum prefixes of a program are
interactions with a memory that has two kinds of actions only, so that the
quantum part of the calculus is confined to one small component.

\paragraph{Mechanisation.} The development is in Rocq, in the directory
\texttt{lqCCS} of the \texttt{testing-theory} library (the earlier
development on quantum distributions is kept apart in \texttt{lqCCS/Prob\_dist}):
\begin{itemize}
\item \texttt{Language/Syntax.v}: syntax with de Bruijn indices for restricted
  channels and bound variables, substitutions, shifts;
\item \texttt{Language/LTS.v}: \cref{fig:prog,fig:mem},
  \cref{def:par}, and the renaming lemmas;
\item \texttt{Language/Congruence.v}: \cref{tab:cgr} and
  \cref{lem:cgr-bisim} (\texttt{cgr\_step\_bisim} for one step, and
  \texttt{cgr\_bisim});
\item \texttt{Harmony\_Lemma/STS\_Harmony.v}: \newline
  \cref{fig:sts} (\texttt{sts});
  \newline \cref{lem:silent} (\texttt{comp\_tau\_inv});
  \cref{lem:shapes}: \newline
  \texttt{TransitionShapeForOutput}, \newline
  \texttt{TransitionShapeForInput}, \newline
  \texttt{TransitionShapeForQuantum}, \newline
  \texttt{TransitionShapeForGuard};
  \cref{lem:red-int}: \newline \texttt{Reduction\_Implies\_Internal};
  \cref{lem:com-shapes}: \texttt{Com\_shapes};
  \cref{lem:int-red}: \newline \texttt{Internal\_Implies\_Reduction}, from
  \texttt{Taus\_Implies\_Reduction} and \texttt{Quantum\_Implies\_Reduction};
  and the main result, \cref{lem:harmony}: \texttt{harmony\_composition};
\item \texttt{Harmony\_Lemma/Memory\_Harmony.v}: the mechanisation also
  defines an LTS directly on the pairs $\conf{\rho}{P}$, by induction on the
  program, and proves that its silent transitions are those of the
  composition $\rho \lpar P$ (\texttt{pre\_harmony}), from which the Harmony
  lemma is restated for this LTS (\texttt{harmony}).
\iffalse%
\item \texttt{Harmony\_Lemma/Memory\_Harmony.v} and \newline
  \texttt{To\_review/Parallel/Parallel\_Harmony.v}: the complements of
  \cref{sec:par}, \cref{rem:prog-level} (\texttt{pre\_harmony}) and
  \cref{lem:par-harmony} (\texttt{parallel\_harmony});
\fi
\ifshowappendix
\item \texttt{Testing\_Pre\_orders/pMust/pMust\_scheduled.v}: distributions,
  the two liftings and the runs, generic; \texttt{pMust\_almost\_sure.v}:
  the almost-sure must, generic; \texttt{pMust\_smyth.v}:
  \cref{def:smyth}, generic, and \texttt{ctx\_pre\_sm\_as};
  \texttt{pMust\_transfer.v}: the transfer of runs along a stuttering
  simulation (\texttt{run\_transfer}, \texttt{pmust\_as\_transfer});
\item \texttt{Testing/Scheduled.v}: the scheduled step of \lqCCS and its instance of the
  generic classes, \texttt{Omega\_bottom}; \texttt{Testing/DistLTS.v}: \cref{tab:dlts}, \cref{def:alt}
  (\texttt{alt\_pre}, \texttt{alt\_as}), the uniform form (\texttt{alt\_pre\_unif},
  \texttt{alt\_pre\_unif\_alt\_pre}), the scheduled-word form (\texttt{follows},
  \texttt{alt\_pre\_det}, \texttt{alt\_pre\_det\_alt\_pre}) and \texttt{drun\_nonzero};
\item \texttt{Testing/Observers.v}: the class of observers, their shape
  (\texttt{menu}, \texttt{S\_T}); \newline
  \texttt{Testing/Smyth\_Soundness.v}: \cref{thm:char}
  (\texttt{alt\_pre\_sound}, \texttt{alt\_pre\_obs\_iff}) and its almost-sure
  counterpart (\texttt{alt\_as\_sound}, \texttt{alt\_as\_obs\_iff}).
\item \texttt{Language/Denotational.v}: \cref{def:den} (\texttt{den},
  \texttt{denot}) and \cref{lem:den-cgr} (\texttt{den\_cgr}, \texttt{denot\_cgr});
\item \texttt{To\_review/Denotational/Denotational\_Harmony.v}: \cref{def:tree}
  (\texttt{sts\_tree}) and \cref{thm:den}: \newline
  \texttt{den\_sound}, \texttt{den\_complete}, \texttt{den\_sts\_iff}.
\fi
\end{itemize}
The quantum model is a parameter: an abstract algebra of super-operators
acting linearly on memories, with the names of the syntax interpreted in it.
